\section{AIME I 1984}
        \begin{enumerate}
        	\item (\textit{AIME I 1984}) Find the value of $a_2+a_4+a_6+a_8+\ldots+a_{98}$ if $a_1$, $a_2$, $a_3\ldots$ is an arithmetic progression with common difference 1, and $a_1+a_2+a_3+\ldots+a_{98}=137$.


 

	\item (\textit{AIME I 1984}) The integer $n$ is the smallest positive multiple of $15$ such that every digit of $n$ is either $8$ or $0$. Compute $\frac{n}{15}$.


 

	\item (\textit{AIME I 1984}) A point $P$ is chosen in the interior of $\triangle ABC$ such that when lines are drawn through $P$ parallel to the sides of $\triangle ABC$, the resulting smaller triangles $t_{1}$, $t_{2}$, and $t_{3}$ in the figure, have areas $4$, $9$, and $49$, respectively. Find the area of $\triangle ABC$.


 
                    \begin{figure}[H]
                        \centering
                        \includegraphics[width=0.5\textwidth]{Problem/AIME_1984_I_3_Problem_Figure_1.png}
                    \end{figure}
                    


 

	\item (\textit{AIME I 1984}) Let $S$ be a list of positive integers--not necessarily distinct--in which the number $68$ appears. The average (arithmetic mean) of the numbers in $S$ is $56$. However, if $68$ is removed, the average of the remaining numbers drops to $55$. What is the largest number that can appear in $S$?


 

	\item (\textit{AIME I 1984}) Determine the value of $ab$ if $\log_8a+\log_4b^2=5$ and $\log_8b+\log_4a^2=7$.


 

	\item (\textit{AIME I 1984}) Three circles, each of radius $3$, are drawn with centers at $(14, 92)$, $(17, 76)$, and $(19, 84)$. A line passing through $(17,76)$ is such that the total area of the parts of the three circles to one side of the line is equal to the total area of the parts of the three circles to the other side of it. What is the absolute value of the slope of this line?


 

	\item (\textit{AIME I 1984}) The function f is defined on the set of integers and satisfies
$f(n)= \begin{cases}  n-3 & \mbox{if }n\ge 1000 \\  f(f(n+5)) & \mbox{if }n<1000 \end{cases}$


 Find $f(84)$.


 

	\item (\textit{AIME I 1984}) The equation $z^6+z^3+1=0$ has complex roots with argument $\theta$ between $90^\circ$ and $180^\circ$ in the complex plane. Determine the degree measure of $\theta$.


 

	\item (\textit{AIME I 1984}) In tetrahedron $ABCD$, edge $AB$ has length 3 cm. The area of face $ABC$ is $15\mbox{cm}^2$ and the area of face $ABD$ is $12 \mbox { cm}^2$. These two faces meet each other at a $30^\circ$ angle. Find the volume of the tetrahedron in $\mbox{cm}^3$.


 

	\item (\textit{AIME I 1984}) Mary told John her score on the American High School Mathematics Examination (AHSME), which was over $80$. From this, John was able to determine the number of problems Mary solved correctly. If Mary's score had been any lower, but still over $80$, John could not have determined this. What was Mary's score? (Recall that the AHSME consists of $30$ multiple choice problems and that one's score, $s$, is computed by the formula $s=30+4c-w$, where $c$ is the number of correct answers and $w$ is the number of wrong answers. Students are not penalized for problems left unanswered.)


 

	\item (\textit{AIME I 1984}) A gardener plants three maple trees, four oaks, and five birch trees in a row. He plants them in random order, each arrangement being equally likely. Let $\frac m n$ in lowest terms be the probability that no two birch trees are next to one another. Find $m+n$.


 

	\item (\textit{AIME I 1984}) A function $f$ is defined for all real numbers and satisfies $f(2+x)=f(2-x)$ and $f(7+x)=f(7-x)$ for all $x$. If $x=0$ is a root for $f(x)=0$, what is the least number of roots $f(x)=0$ must have in the interval $-1000\leq x \leq 1000$?


 

	\item (\textit{AIME I 1984}) Find the value of $10\cot(\cot^{-1}3+\cot^{-1}7+\cot^{-1}13+\cot^{-1}21).$


 

	\item (\textit{AIME I 1984}) What is the largest even integer that cannot be written as the sum of two odd composite numbers?


 

	\item (\textit{AIME I 1984}) Determine $w^2+x^2+y^2+z^2$ if


 $\frac{x^2}{2^2-1}+\frac{y^2}{2^2-3^2}+\frac{z^2}{2^2-5^2}+\frac{w^2}{2^2-7^2}=1$$\frac{x^2}{4^2-1}+\frac{y^2}{4^2-3^2}+\frac{z^2}{4^2-5^2}+\frac{w^2}{4^2-7^2}=1$$\frac{x^2}{6^2-1}+\frac{y^2}{6^2-3^2}+\frac{z^2}{6^2-5^2}+\frac{w^2}{6^2-7^2}=1$$\frac{x^2}{8^2-1}+\frac{y^2}{8^2-3^2}+\frac{z^2}{8^2-5^2}+\frac{w^2}{8^2-7^2}=1$

\end{enumerate}